\subsubsection{Examples}
\label{subsubsec:modular-examples}

The following examples record exact consequences of
\cref{prop:modular-formula}.
\par\medskip

\begin{example}[The projective planes]
\label{ex:projective-planes}
Let $e^2=1$ generate $H_2(\mathbb{CP}^2;\mathbb Z)$ and let
$\bar e^2=-1$ generate $H_2(\overline{\mathbb{CP}}^{\,2};\mathbb Z)$.  For
$n\in\mathbb Z$,
\[
  \widehat{\CY}_{\cC,\eta}(\mathbb{CP}^2,ne)
    =\tau_+(\cC)\theta_{g_\eta}^{-n^2},
  \qquad
  \widehat{\CY}_{\cC,\eta}(\overline{\mathbb{CP}}^{\,2},n\bar e)
    =\tau_-(\cC)\theta_{g_\eta}^{n^2}.
\]
Thus the entire marked function on either projective plane is determined by
one Gauss sum and the twist of $g_\eta$.
\end{example}

\begin{example}[Diagonal rational surfaces]
\label{ex:diagonal-rational-surfaces}
For
\[
 X=p\mathbb{CP}^2\#q\overline{\mathbb{CP}}^{\,2},
 \qquad
 \alpha=(a_1,\ldots,a_p;b_1,\ldots,b_q),
\]
the modular closed formula gives
\[
 \widehat{\CY}_{\cC,\eta}(X,\alpha)
 =\tau_+(\cC)^p\tau_-(\cC)^q
  \theta_{g_\eta}^{-
    (\sum_r a_r^2-\sum_s b_s^2)}.
\]
In particular, for $\mathbb{CP}^2\#2\overline{\mathbb{CP}}^{\,2}$ the
exponent is $-(a^2-b^2-c^2)$.
\end{example}

\begin{example}[The hyperbolic block]
\label{ex:s2xs2}
Write $H=S^2\times S^2$ and choose the hyperbolic basis $u,v$, with
$u^2=v^2=0$ and $u\cdot v=1$.  Then
\[
  \widehat{\CY}_{\cC,\eta}(H,au+bv)
  =\tau_+(\cC)\tau_-(\cC)\theta_{g_\eta}^{-2ab}.
\]
For the zero marking $a=b=0$, this is precisely the factor introduced by a
hyperbolic stabilization and cancelled in the stabilization argument.
\end{example}

\begin{example}[$K3$]
\label{ex:k3-modular}
With the complex orientation, $b_2^+(K3)=3$ and $b_2^-(K3)=19$; hence
\[
  \widehat{\CY}_{\cC,\eta}(K3,\alpha)
  =\tau_+(\cC)^3\tau_-(\cC)^{19}
   \theta_{g_\eta}^{-Q_{K3}(\alpha,\alpha)}.
\]
No Kirby-diagram contraction is needed once the intersection form and the
self-intersection of $\alpha$ are known.
\end{example}

Independent numerical checks of these formulas, together with the
runtime, cloud-cost, classical-scaling, and quantum-query estimates for the
direct evaluator, are collected in
Appendix~\ref{app:computational-record}.
